\begin{exercise}{58}
Show the following assertions:
\begin{enumerate}[label=(\arabic*)]
\item The map $\varnothing\to\{*\}$ in $\mathrm{Set}$ generates the set of injections. What is $\chi_R(\varnothing\to *)$? Spell out
the factorization obtained by the small object argument for a general map $f\colon M\to N$
of sets.
\item The map $\{*,*\}\to\{*\}$ generates the class of surjections. What is $\chi_R(\{*,*\}\to *)$? Spell
out the factorization obtained by the small object argument for a general map $f\colon M\to
N$ of sets.
\end{enumerate}

\begin{proof}
Write $i\colon\varnothing\to\{*\}$ and
$c\colon\{*,*\}\to\{*\}$.
A lifting square from $i$ to $p$ specifies a point of the target
of $p$; a diagonal is a preimage of that point.
A square from $c$ specifies two points with the same image
under $p$; a diagonal exists exactly when those points coincide.
Thus $\chi_R(i)$ is the class of surjections and
$\chi_R(c)$ is the class of injections.

An injection has the LLP against every surjection by choosing
preimages off its image; a non-injection fails this property
against $c$, after separating two points that it identifies.
A surjection has the LLP against every injection because
commutativity makes the top map constant on its fibers;
a non-surjection fails against the inclusion of its image.
Consequently $\chi(i)$ consists of the injections and
$\chi(c)$ of the surjections.
These classes are saturated by Lemma~1.3.8.
Conversely, every injection is a pushout of a coproduct of $i$
(one copy per point of its complement), and every surjection
is a pushout of a coproduct of $c$ (one copy per kernel pair).
This proves both generation claims.

To compute the factorizations, use the construction of
Proposition~1.3.9, attaching a cell for every commutative square
at every stage, starting with $E^0=M$.

(1) For $i$, at every stage the squares are indexed by $N$.
Thus $E^{k+1}=E^k\amalg N$, with the new copy mapping identically
to $N$. The resulting factorization is
\[
 M\lhook\joinrel\longrightarrow
 M\amalg\bigl(N\times\mathbb N_{>0}\bigr)
 \xrightarrow{\ p\ }N,
 \qquad p(m)=f(m),\quad p(n,r)=n.
\]
Indeed, $E^k=M\amalg(N\times\{1,\ldots,k\})$, and the displayed
middle object is its sequential colimit.
The first map is injective and $p$ is surjective.

(2) For $c$, the first stage is indexed by all pairs
$(a,b)\in M^2$ with $f(a)=f(b)$, including diagonal pairs.
The pushout identifies exactly these pairs, giving
\[
 M\twoheadrightarrow M/R\xrightarrow{\ \bar f\ }N,
 \qquad a\mathrel R b\ \Longleftrightarrow\ f(a)=f(b),
 \qquad \bar f([a])=f(a).
\]
The map $\bar f$ is injective.
At every later stage the squares correspond to diagonal pairs
in $M/R$. Their pushouts identify each point with itself,
so the transitions are isomorphisms and
$E^{\omega}\cong M/R$.
This is the quotient surjection followed by the induced injection.
\end{proof}
\end{exercise}
